\chapter{Rassegna della letteratura}
\label{ch:letteratura}

Questo capitolo colloca il lavoro all'interno di cinque filoni di ricerca.
Il primo definisce il premio per il rischio di volatilità nei mercati
tradizionali e ne discute le spiegazioni teoriche
(\S\ref{sec:lit-vrp-tradizionale}). Il secondo esamina le strategie
sistematiche di vendita di volatilità e il loro profilo di rischio
(\S\ref{sec:lit-short-vol}). Il terzo passa al mercato delle opzioni su
criptovalute, dove il fenomeno assume dimensioni particolari
(\S\ref{sec:lit-crypto-options}). Il quarto e il quinto trattano gli
strumenti metodologici impiegati nella tesi: la stima della volatilità
realizzata (\S\ref{sec:lit-realized-vol}) e l'estrazione della densità
risk-neutral (\S\ref{sec:lit-rnd}), oltre ai criteri di dimensionamento del
rischio (\S\ref{sec:lit-sizing}). Il capitolo si chiude posizionando il
contributo di questa tesi (\S\ref{sec:lit-posizionamento}).

\section{Il Volatility Risk Premium nei mercati tradizionali}
\label{sec:lit-vrp-tradizionale}

Il premio per il rischio di volatilità (VRP) è la differenza sistematica
tra la volatilità implicita nei prezzi delle opzioni e la volatilità che il
sottostante realizza successivamente. Empiricamente questa differenza è
positiva e persistente: chi vende opzioni, e dunque vende volatilità,
incassa in media un premio, mentre chi le acquista paga un costo atteso
superiore al valore attuariale della protezione ottenuta.

La prima evidenza sistematica del fenomeno su singole azioni e indici si
deve a \citet{bakshi2003delta}, che studiano i rendimenti di posizioni in
opzioni coperte in delta (\emph{delta-hedged gains}). Se il modello di
Black--Scholes valesse esattamente e non esistesse alcun premio per il
rischio di volatilità, il guadagno di una posizione coperta in delta
dovrebbe essere nullo in media. Gli autori documentano invece guadagni
sistematicamente negativi per chi detiene le opzioni, e li interpretano
come prezzo di mercato di un rischio di volatilità priced negativamente:
gli investitori sono disposti a pagare per detenere strumenti che
guadagnano quando la volatilità sale.

\citet{carr2009variance} formalizzano il concetto attraverso i
\emph{variance swap}, la cui gamba fissa (il tasso di variance swap) è
replicabile in modo model-free da un portafoglio statico di opzioni su un
continuum di strike. La differenza tra il tasso di variance swap
(volatilità attesa sotto la misura risk-neutral $\mathbb{Q}$) e la varianza
effettivamente realizzata (sotto la misura fisica $\mathbb{P}$) misura
direttamente il variance risk premium. Su indici e single name statunitensi
il premio è ampiamente negativo per i compratori, coerentemente con
\citet{bakshi2003delta}.

\citet{bollerslev2009expected} mostrano che il VRP non è solo un premio
statico ma ha contenuto informativo: la differenza tra varianza implicita e
varianza realizzata prevede i rendimenti azionari futuri, con potere
predittivo particolarmente forte sull'orizzonte trimestrale. Gli autori
collegano il premio all'avversione al rischio aggregata e alla variabilità
temporale dell'incertezza economica, offrendo un fondamento di equilibrio
al fenomeno. Le principali spiegazioni teoriche del VRP, non
mutuamente esclusive, sono infatti tre: l'avversione degli investitori a
uno stato del mondo (alta volatilità) in cui la ricchezza marginale è
elevata; la domanda strutturale di copertura, che spinge i prezzi delle
opzioni sopra il loro valore attuariale; e i vincoli di capitale e di
inventario dei market maker, che richiedono un compenso per assorbire
l'offerta di rischio di volatilità.

\section{Strategie sistematiche di vendita di volatilità}
\label{sec:lit-short-vol}

Se il VRP è un premio persistente, la sua cattura sistematica è in linea di
principio possibile vendendo opzioni. La letteratura sulle strategie
passive di vendita di volatilità nasce con \citet{whaley2002return}, che
analizza l'indice CBOE S\&P~500 BuyWrite (BXM), una strategia
\emph{buy-write} che detiene l'indice e vende sistematicamente call
at-the-money. Whaley documenta, sul periodo 1988--2001, rendimenti
corretti per il rischio superiori a quelli del sottostante e, soprattutto,
una marcata riduzione della volatilità del portafoglio: la vendita
sistematica della call converte parte del premio per il rischio di
volatilità in reddito, al costo di rinunciare all'upside.

Questa letteratura evidenzia però anche il profilo di rischio asimmetrico
delle strategie short volatility. Il payoff tipico è composto da molti
piccoli guadagni intervallati da rare perdite di grande entità, un profilo
sintetizzato dalla metafora del \emph{picking up nickels in front of a
steamroller}. Tre implicazioni sono rilevanti per questa tesi. Primo, le
metriche basate sui soli momenti secondi (come lo Sharpe ratio) possono
sovrastimare la qualità di queste strategie, che concentrano il rischio
nelle code. Secondo, il dimensionamento della posizione e la gestione del
rischio di coda sono determinanti quanto la selezione dell'opzione. Terzo,
l'uso di stop loss su strategie short volatility è problematico, perché la
componente di theta decay che genera il profitto richiede tempo e viene
troncata da uscite premature, un punto che la presente tesi verifica
empiricamente nel Capitolo~\ref{ch:percorso}.

\section{Il mercato delle opzioni su criptovalute}
\label{sec:lit-crypto-options}

Il mercato delle opzioni su Bitcoin è giovane ma già oggetto di una
letteratura empirica in rapida crescita. La quasi totalità dell'attività si
concentra su Deribit, che controlla oltre l'80\% del volume globale di
opzioni crypto \citep{alexander2023inverse}. Una specificità tecnica del
mercato è che gli strumenti quotati sono \emph{inverse options}: premio e
regolamento sono denominati in Bitcoin, mentre lo strike è in dollari.
\citet{alexander2023inverse} mostrano che l'applicazione ingenua della
formula di Black--Scholes standard alle inverse options produce volatilità
implicite distorte, e che la volatilità implicita corretta per la struttura
inverse è quasi sempre superiore a quella ottenuta con la formula standard:
un dettaglio metodologico di cui la costruzione della volatility
surface deve tenere conto.

Il contributo più direttamente rilevante per questa tesi è quello di
\citet{alexander2020bitcoin}, che costruiscono un indice di volatilità
implicita per Bitcoin analogo al VIX, il precursore di quello che
Deribit avrebbe poi standardizzato nell'indice DVOL, a partire da oltre
sette milioni di prezzi di opzioni campionati ogni 15 minuti. Gli autori
documentano un premio per il rischio di varianza su Bitcoin
significativamente più ampio di quello dei mercati azionari. Studi
successivi ne quantificano la magnitudine: \citet{almeida2024risk} stimano
un variance risk premium annualizzato su Bitcoin compreso tra il 7\% e il
14\%, contro il 2\% circa dell'S\&P~500, un premio dunque di un ordine
di grandezza superiore. Un risultato particolarmente rilevante del loro
lavoro è che il VRP è \emph{più elevato nei regimi di bassa volatilità}
(0,10--0,17) che in quelli di alta volatilità (0,04--0,12), a indicare che
la preoccupazione degli investitori per il rischio di varianza non è
semplicemente proporzionale al livello corrente di volatilità.

Sul segno e sulla catturabilità del premio, \citet{winkel2023pricing}
stimano i pricing kernel impliciti nelle opzioni Deribit, ossia il rapporto
tra densità risk-neutral e densità fisica, e trovano che gli investitori
pagano un premio sostanziale per assicurarsi contro i movimenti di prezzo
di breve termine, premio decrescente con la scadenza. La letteratura
converge sul fatto che il premio per il rischio di volatilità osservato su
Deribit è negativo e significativo, cosicché le strategie di vendita di
volatilità sono attese generare performance corrette per il rischio
positive. È precisamente questa l'ipotesi che la presente tesi mette alla
prova in modo sistematico e con costi di transazione realistici.

\section{Stima della volatilità realizzata}
\label{sec:lit-realized-vol}

La misurazione della volatilità realizzata è centrale per quantificare il
VRP, poiché costituisce il termine di confronto della volatilità implicita.
Lo stimatore close-to-close, basato sui soli prezzi di chiusura, ignora
l'informazione contenuta nell'escursione infragiornaliera dei prezzi.
\citet{parkinson1980extreme} propone uno stimatore basato sul range
high--low che, sotto l'ipotesi di moto browniano senza drift, è circa
cinque volte più efficiente dello stimatore close-to-close. Per Bitcoin,
che è negoziato senza interruzione 24 ore su 24 e 7 giorni su 7, lo
stimatore di Parkinson è particolarmente appropriato: non esistono
discontinuità di chiusura e l'informazione infragiornaliera è pienamente
disponibile e continua. Per questa ragione la tesi adotta la volatilità
realizzata di Parkinson a 30 giorni come misura di riferimento.

Sul fronte del \emph{forecasting} della volatilità, non impiegato nella
strategia finale ma rilevante per gli sviluppi futuri, i due riferimenti
canonici sono il modello GARCH di \citet{bollerslev1986generalized}, che
cattura la persistenza (clustering) della volatilità condizionata, e il
modello HAR-RV di \citet{corsi2009simple}, che approssima la memoria lunga
della volatilità realizzata combinando componenti giornaliere, settimanali
e mensili. Entrambi sono discussi nel Capitolo~\ref{ch:limitazioni} come
possibili vie per riabilitare un segnale di \emph{timing} del premio.

\section{Estrazione della densità risk-neutral}
\label{sec:lit-rnd}

Il risultato fondamentale sull'estrazione della densità risk-neutral (RND)
implicita nei prezzi delle opzioni è quello di \citet{breeden1978prices}:
la densità risk-neutral del sottostante alla scadenza è proporzionale alla
derivata seconda del prezzo della call rispetto allo strike. In pratica,
disponendo di prezzi di opzioni su un insieme discreto di strike, la RND si
ottiene interpolando la superficie di volatilità implicita (tipicamente con
spline o parametrizzazioni tipo SVI), convertendo in prezzi di call
continui e differenziando numericamente due volte.

Questa infrastruttura è stata impiegata nella presente ricerca per
costruire un segnale di selezione delle opzioni, confrontando la RND (sotto
la misura $\mathbb{Q}$) con una stima della densità fisica (sotto la misura
$\mathbb{P}$) ottenuta via kernel density estimation. La letteratura
avverte tuttavia sulla fragilità del confronto \emph{pointwise} tra due
densità entrambe stimate con errore: la RND è sensibile agli errori di
interpolazione sulle code, e la densità fisica dipende in modo critico da
assunzioni di stazionarietà e dalla scelta della finestra. Come documentato
nel Capitolo~\ref{ch:percorso}, questa fragilità è la ragione empirica per
cui il segnale basato sul gap RND--KDE si è rivelato privo di potere
predittivo, pur restando la RND un utile strumento di contesto. La stessa
metodologia di stima dei pricing kernel è al centro del lavoro di
\citet{winkel2023pricing} e \citet{almeida2024risk} sul mercato Bitcoin.

\section{Criteri di dimensionamento e gestione del rischio}
\label{sec:lit-sizing}

Il problema del dimensionamento ottimo della posizione ha il suo
riferimento classico nel criterio di \citet{kelly1956new}, che massimizza
il tasso di crescita atteso del capitale nel lungo periodo. Il criterio di
Kelly è tuttavia noto per essere aggressivo e per essere estremamente
sensibile agli errori di stima dei parametri (probabilità di vincita e
rapporto payoff): quando questi sono stimati su campioni ridotti,
l'incertezza sulla frazione ottima è ampia e l'applicazione della frazione
piena espone a drawdown severi. La pratica di mercato adotta quindi
frazioni ridotte del Kelly (\emph{fractional Kelly}, tipicamente tra un
mezzo e un decimo), che sacrificano una piccola parte del tasso di crescita
atteso in cambio di una netta riduzione della variabilità del capitale.
Nella presente tesi il Half-Kelly è utilizzato non come metodo di
dimensionamento ma come \emph{sanity check} della ragionevolezza delle
configurazioni adottate (Capitolo~\ref{ch:money-management}), coerentemente
con l'osservazione che il vincolo operativo rilevante per una strategia
short volatility non è tanto il rischio di rovina quanto la tolleranza al
drawdown.

\section{Posizionamento di questa tesi}
\label{sec:lit-posizionamento}

Rispetto alla letteratura descritta, il contributo di questa tesi si
articola su tre piani. Primo, sul piano \emph{applicativo}: mentre gran
parte degli studi sul VRP crypto si concentra sulla misurazione del premio
(la sua esistenza, la sua magnitudine, il suo segno), questa tesi affronta
la domanda operativa di se e come il premio sia catturabile da una
strategia replicabile, al netto di costi di transazione realistici,
dell'esecuzione T+1 e della copertura delta. Secondo, sul piano
\emph{metodologico}: il lavoro documenta in modo trasparente non solo la
strategia sopravvissuta, ma anche le cinque ipotesi testate e scartate:
un antidoto al \emph{publication bias} che affligge la letteratura sulle
strategie di trading, dove tipicamente vengono riportati solo i risultati
positivi. In particolare, l'individuazione dell'artefatto da opzioni
illiquide con bid nullo (Capitolo~\ref{ch:percorso}) è un monito
metodologico di portata generale. Terzo, sul piano della \emph{robustezza}:
la validazione cross-asset su opzioni SPY, con la stessa strategia e senza
alcuna riottimizzazione, funge da test out-of-sample su una asset class
diversa e su un orizzonte temporale molto più lungo, isolando il meccanismo
di cattura del premio dalle specificità del campione Bitcoin.
